% Abstract of "Communication, Coherence and Collective Agents".
% src: framing pass; hyp:int-h1 to hyp:int-h5 and their status sentences in sec:int-hypotheses; sec:dis-status; introduction of sec:opn with tab:opn-status and tab:opn-status-b; findings frame-03, frame-06, pos-03, pos-08, frm-09 of frame/00/items_1 to items_3 (honored by the narrower opening, the sentence on the meta-agent definition, the declared predictive objective, and by quoting no replicate count)
\begin{abstract}
We develop a linear-Gaussian active-inference framework for deciding when communicating agents count as a collective agent. Agents send coded beliefs about an external world over noisy channels; receivers lift them by prior completion. Five hypotheses organize the work. H1, collectives as graded meta-agents: the boundary clause, an exact path blanket, is proved and selects nothing, since it holds for every candidate with separate, independently noisy channels; collectives are selected by definition, through contrast against a declared comparison medium; partial results on closure, interaction necessity and agency are proved under stated conditions, but no analyzed population is shown to meet every clause of the meta-agent definition, whose joint satisfiability is open. H2, that literal receivers misread the joint structure of messages, is proved in stated scopes, with undervaluation elsewhere a numerical observation. H3, formation through coherent message passing, is proved only in a two-group model and, at finite policy precision above a threshold, only as reachability in a frozen-calibration mean-field map, with the stochastic step from the full model open; from a start exchangeable given the world path the split fails the translation-aware exact-fusion persistence clause, and formation beyond that model is a conjecture. H4, growth of informational correlation lengths and times with scale, is proved under a moment bound and a named contraction, largely as upper bounds. H5, meta-agents as agents with an abstract physics, joined across levels by top-down priors, bottom-up evidence and renormalization, is proved, with named exceptions, under stated conditions for declared linear coarse states; its exact conditional theorem selects no meta-agent, its approximate form is proved only for stable linear-Gaussian collectives whose members run the consensus filter and is open beyond them, its sensing results concern a declared predictive objective that the one-step expected free energy does not imply, and formation by conformity is a conjecture.

We address ten open problems posed in the literature on collective active inference, from collective agency to empirical identification. None is closed in the generality posed, and energetic viability and evolutionary individuality are addressed only by scoped models. Proofs, numerical checks, secondary results and eight extensions addressing these problems are in a supplement.
\end{abstract}
